% Generated by tools/generate_statements_only_paper.py.
% Do not edit this derived file; edit the canonical paper source instead.

\subsection{The finite caliper criterion}
\label{subsec:finite-caliper}

We record the finite caliper criterion used below; its verification is omitted
from this reading edition.

Let $\mathsf R$ and $\mathsf J$ denote counterclockwise rotations through $2\pi/3$ and
$\pi/2$, respectively.  For a nonempty finite set $S\subset\mathbb R^2$ and
a vector $N$, put
\begin{equation}
 W_S(N)=\sum_{k=0}^2\max_{z\in S}\langle z,\mathsf R^kN\rangle .
 \label{eq:caliper-support-sum}
\end{equation}

\begin{theorem}[Finite caliper certificate]
\label{thm:cert-caliper}
If $S$ has at least two distinct points, then $S$ is contained in a closed
unit equilateral triangle if and only if there are distinct $p,q\in S$ and
$\varepsilon\in\{-1,1\}$ for which
\begin{equation}
 W_S\bigl(\varepsilon\mathsf J(q-p)\bigr)
 \le \frac{\sqrt3}{2}\lVert p-q\rVert .
 \label{eq:finite-caliper-test}
\end{equation}
If $S$ consists of one point, it is contained in equilateral triangles of
arbitrarily small positive side length.

For $S_x=\{O,A,B,x\}$, every clause of
\eqref{eq:finite-caliper-test} is a finite polynomial clause of degree at
most four in the coordinates of $A,B,x$; at most $12\cdot4^3$ clauses are
required, and repeated point labels merely delete zero-length pairs.
\end{theorem}

% Proof environment omitted (source: appendix_certificates.tex).
